\chapter{Methodology}
\label{ch:methodology}

\section{Research Design and Overall Approach}
\label{sec:research_design}

This thesis adopts a constructive research design, combining framework-driven assessment with empirical tool development. Rather than observing cloud sovereignty as a static phenomenon, the research operationalises it through two complementary instruments that can be applied, reproduced, and extended by practitioners. The first instrument — the Sovereignty Scoring Framework (SovScore) — measures the degree of digital sovereignty a cloud provider offers to EU customers. The second — the Cloud Readiness Assessment — measures how prepared each provider is to substitute AWS-equivalent managed services for a European enterprise workload. Both instruments are grounded in publicly available regulatory documents and evaluated against a curated set of 18 European cloud providers.

The two experiments are described in detail in Sections~\ref{sec:sovscore_methodology} and~\ref{sec:readiness_methodology} respectively. Both were conducted using an open-source web tool developed as part of this research, described in the following section.


% ─────────────────────────────────────────────────────────────────────────────
\section{CloudSov Tool: Implementation and Usage}
\label{sec:cloudsov_tool}
% ─────────────────────────────────────────────────────────────────────────────

\subsection{Overview and Availability}

To ensure reproducibility and accessibility, both experiments were operationalised within a dedicated open-source web application named \textbf{CloudSov}. The tool is publicly available on GitHub and can be run locally without a cloud account or paid subscription. It provides a graphical interface for conducting the SovScore assessment and the Cloud Readiness Assessment in a structured, comparable manner across all 18 EU providers.

\begin{quote}
\textit{Repository:} \texttt{https://github.com/aaeu-ihono/cloudsov-tool.git}
\textit{Licence:} MIT
\end{quote}

\noindent The tool is built with a Python/FastAPI backend and a React~19 (Vite) frontend, communicating over a local REST~API. Provider responses are stored as structured JSON files that can be audited independently of the interface.

\subsection{Getting Started}

Prerequisites: Python~3.10+, Node.js~18+, and \texttt{npm}. To launch the tool:

\begin{enumerate}
  \item Clone the repository: \texttt{git clone https://github.com/aaeu-ihono/cloudsov-tool.git}
  \item Start the backend: \texttt{cd backend \&\& pip install -r requirements.txt \&\& uvicorn main:app --reload}
  \item Start the frontend: \texttt{cd frontend \&\& npm install \&\& npm run dev}
  \item Open a browser at \texttt{http://localhost:5173}
\end{enumerate}

No environment variables or external API keys are required. All provider data is bundled with the repository under the \texttt{survey/} and \texttt{readiness/} directories.

\subsection{Interface Walkthrough}

The tool interface is divided into two primary modules, accessible from the top navigation bar: \textbf{SovScore} and \textbf{Readiness Assessment}. Figures~\ref{fig:tool_sovscore_grid}–\ref{fig:tool_readiness_gap} illustrate each section. The numbered annotations in the figures correspond to the legend in Table~\ref{tab:tool_legend}.

\begin{table}[htbp]
\centering
\caption{Interface annotation legend — numbered regions shown in Figures~\ref{fig:tool_sovscore_grid}–\ref{fig:tool_readiness_gap}}
\label{tab:tool_legend}
\begin{tabularx}{\linewidth}{cX}
\toprule
\textbf{No.} & \textbf{Description} \\
\midrule
\circled{1} & Provider selector and add/remove controls — manage which providers appear in the assessment grid \\
\circled{2} & SOV objective column headers — each column represents one of the eight sovereignty objectives (SOV-1 to SOV-8) \\
\circled{3} & Verdict input cells — each cell accepts Y (Yes), P (Partial), or N (No) for the corresponding provider–question pair \\
\circled{4} & Computed SovScore column — displays the weighted aggregate score (0–100) updated in real time \\
\circled{5} & SEAL badge — colour-coded maturity level (SEAL~0 to SEAL~4) derived from the score \\
\circled{6} & Dimension breakdown radar/bar chart — shows per-dimension readiness scores for a selected provider \\
\circled{7} & Provider toggle panel — enables or disables individual providers on the readiness line chart \\
\circled{8} & Gap table — lists all 60 services ranked by average gap from AWS baseline; expandable rows show per-provider status \\
\bottomrule
\end{tabularx}
\end{table}

\noindent To use the \textbf{SovScore} module (Figure~\ref{fig:tool_sovscore_grid}): select the providers to compare using the panel at the top \circled{1}, then work through each column \circled{2}, entering verdicts in the cells \circled{3}. The score \circled{4} and SEAL badge \circled{5} update automatically. Hovering over a cell reveals the full question text and an evidence note field where source URLs can be recorded.

\noindent To use the \textbf{Readiness Assessment} module (Figure~\ref{fig:tool_readiness_chart}): toggle providers on or off using the pill buttons \circled{7}. The line chart displays dimension scores as a time-series of 13 capability dimensions. Hovering over any data point reveals a tooltip listing which of the 60 services are available (Y), partial (P), or missing (N) for that provider in that dimension. The gap table \circled{8} (Figure~\ref{fig:tool_readiness_gap}) ranks services by average deficit and allows drill-down per provider.

\begin{figure}[htbp]
  \centering
  \fbox{\begin{minipage}{0.92\linewidth}\centering\vspace{4cm}
    \textit{[Screenshot: SovScore assessment grid — annotated with numbers \circled{1}–\circled{5}]}
  \vspace{4cm}\end{minipage}}
  \caption{CloudSov SovScore module: provider grid with verdict inputs and computed scores.}
  \label{fig:tool_sovscore_grid}
\end{figure}

\begin{figure}[htbp]
  \centering
  \fbox{\begin{minipage}{0.92\linewidth}\centering\vspace{4cm}
    \textit{[Screenshot: Readiness Assessment line chart — annotated with \circled{6} and \circled{7}]}
  \vspace{4cm}\end{minipage}}
  \caption{CloudSov Readiness Assessment module: dimension scores across providers.}
  \label{fig:tool_readiness_chart}
\end{figure}

\begin{figure}[htbp]
  \centering
  \fbox{\begin{minipage}{0.92\linewidth}\centering\vspace{4cm}
    \textit{[Screenshot: Gap table with expanded service row — annotated with \circled{8}]}
  \vspace{4cm}\end{minipage}}
  \caption{CloudSov gap table: services ranked by average deficit from the AWS baseline.}
  \label{fig:tool_readiness_gap}
\end{figure}

\subsection{Advantages, Limitations, and Controversies}

\paragraph{Advantages.}
CloudSov addresses the absence of a unified, comparable, open evaluation framework for EU cloud sovereignty. By combining two complementary instruments in a single interface, it enables side-by-side comparison across providers that would otherwise require bespoke spreadsheets and manual aggregation. The tool is technology-neutral, provider-agnostic, and entirely reproducible: any researcher with access to the repository can re-run or extend the assessment.

\paragraph{Limitations.}
The assessments rely on publicly available documentation and self-reported provider claims. No independent technical audit of the verdicts was conducted. Providers with sparse public documentation may receive lower scores not because they fail to meet the criteria, but because sufficient evidence could not be established — a documentation gap rather than a capability gap. Furthermore, the SovScore weights are prescribed by the EC~DG~DIGIT framework and were not derived empirically; alternative weighting schemes would yield different rankings.

\paragraph{Controversies.}
The use of AWS as a readiness baseline has inherent tensions: it implicitly treats the US hyperscaler's service model as the normative standard against which EU alternatives are measured, which may undervalue architectural choices made by EU providers (e.g., OpenStack-based deployments that prioritise interoperability over proprietary services). Additionally, the binary gap model (available / partial / missing) cannot capture quality, performance, or SLA differences between an EU provider's service and its AWS equivalent. These limitations are acknowledged in the discussion of findings in Chapter~8.

% ─────────────────────────────────────────────────────────────────────────────
\section{Sovereignty Scoring Framework (SovScore)}
\label{sec:sovscore_methodology}
% ─────────────────────────────────────────────────────────────────────────────

\subsection{Theoretical Basis and Framework Selection}

Cloud computing has historically been evaluated along performance, reliability, and cost dimensions~\cite{armbrust2010view, NIST_SP800145_2011}. The policy dimension — concerning legal jurisdiction, data control, and strategic independence from non-EU actors — has remained largely qualitative and inconsistently defined across the literature. This experiment addresses that gap by grounding the sovereignty assessment in the most authoritative publicly available specification for EU cloud governance: the European Commission Directorate-General for Digital Services \textit{Cloud Sovereignty Framework v1.2.1}, published in October~2025~\cite{EC_DIGIT_CloudSovereignty_2025}.

The decision to adopt the EC~DG~DIGIT framework over alternative sovereign cloud criteria%
\footnote{Alternative frameworks considered include the ENISA EUCS certification scheme~\cite{ENISA_EUCS_2023} and the GAIA-X Trust Framework~\cite{GAIAX_Architecture_2023}. The EC~DG~DIGIT document was preferred because it provides explicit, question-level evaluation criteria that map directly to measurable provider attributes, rather than certification labels that presuppose a third-party audit process.}
was motivated by three factors. First, it is produced by the Commission body responsible for EU digital policy, giving it direct regulatory standing. Second, it decomposes sovereignty into eight thematic objectives with explicit sub-criteria and a five-level maturity scale. Third, it is technology-neutral and provider-agnostic, making it applicable across the heterogeneous landscape of European cloud providers without presupposing any particular architectural pattern.

\subsection{SOV Objectives and Weighting Scheme}
\label{subsec:sov_objectives}

The framework defines eight \textit{SOV objectives}, each representing a distinct dimension of digital sovereignty. Each objective is assigned a percentage weight reflecting its relative importance to overall sovereignty, as specified in the EC~DG~DIGIT framework. These weights sum to 100\% and are held constant across all provider assessments to ensure comparability.

Table~\ref{tab:sov_objectives} summarises the eight objectives with their weights and minimum SEAL thresholds. Tables~\ref{tab:sov_objectives_a} and~\ref{tab:sov_objectives_b} list the full evaluation questions for each objective.

\begin{table}[htbp]
\centering
\caption{SOV objectives, weights, question counts, and minimum SEAL thresholds}
\label{tab:sov_objectives}
\begin{tabular}{llrrr}
\toprule
\textbf{ID} & \textbf{Objective} & \textbf{Weight (\%)} & \textbf{Questions} & \textbf{Min.\ SEAL} \\
\midrule
SOV-1 & Strategic Sovereignty      & 15 & 6 & 3 \\
SOV-2 & Legal \& Jurisdictional    & 10 & 5 & 4 \\
SOV-3 & Data \& AI Sovereignty     & 10 & 5 & 3 \\
SOV-4 & Operational Sovereignty    & 15 & 6 & 2 \\
SOV-5 & Supply Chain               & 20 & 5 & 2 \\
SOV-6 & Technology                 & 15 & 4 & 2 \\
SOV-7 & Security \& Compliance     & 10 & 6 & 3 \\
SOV-8 & Environmental              &  5 & 4 & 1 \\
\midrule
       & \textbf{Total}             & \textbf{100} & \textbf{41} & \\
\bottomrule
\end{tabular}
\end{table}

% ── SOV-1 to SOV-4 ───────────────────────────────────────────────────────────
\begin{table}[htbp]
\centering
\footnotesize
\renewcommand{\arraystretch}{1.25}
\caption{SOV objectives SOV-1 to SOV-4: weights, min.\ SEAL thresholds, and evaluation questions}
\label{tab:sov_objectives_a}
\begin{tabularx}{\linewidth}{|c|X|X|X|X|}
\hline
 &
\textbf{SOV-1}\newline Strategic Sovereignty\newline \textit{15\% · min~SEAL-3} &
\textbf{SOV-2}\newline Legal \& Jurisdictional\newline \textit{10\% · min~SEAL-4} &
\textbf{SOV-3}\newline Data \& AI Sovereignty\newline \textit{10\% · min~SEAL-3} &
\textbf{SOV-4}\newline Operational Sovereignty\newline \textit{15\% · min~SEAL-2} \\
\hline\hline
Q1 &
Decisive authority bodies within EU? &
Governed exclusively by EU member law? &
Customer-managed encryption keys (CMEK)? &
Workload portability without vendor lock-in? \\
\hline
Q2 &
Structural protection against non-EU takeover? &
Free from CLOUD Act and extraterritorial laws? &
Full auditability of data and AI access? &
EU-managed operations, no non-EU vendor dependency? \\
\hline
Q3 &
Primarily EU-sourced financing? &
No channel for non-EU compelled data access? &
Verifiable irreversible data deletion? &
EU-based talent pool for operations? \\
\hline
Q4 &
Significant EU investment and job creation? &
Free from non-EU export control regimes? &
Data storage and processing EU-only? &
All operational support from within EU? \\
\hline
Q5 &
Active in GAIA-X and EU sovereignty initiatives? &
IP registered and developed within EU? &
AI models trained under EU control? &
Technical documentation and source code available? \\
\hline
Q6 &
Operational continuity if non-EU support withdrawn? &
--- &
--- &
All critical suppliers under EU jurisdiction? \\
\hline
\end{tabularx}
\end{table}

% ── SOV-5 to SOV-8 ───────────────────────────────────────────────────────────
\begin{table}[htbp]
\centering
\footnotesize
\renewcommand{\arraystretch}{1.25}
\caption{SOV objectives SOV-5 to SOV-8: weights, min.\ SEAL thresholds, and evaluation questions}
\label{tab:sov_objectives_b}
\begin{tabularx}{\linewidth}{|c|X|X|X|X|}
\hline
 &
\textbf{SOV-5}\newline Supply Chain\newline \textit{20\% · min~SEAL-2} &
\textbf{SOV-6}\newline Technology\newline \textit{15\% · min~SEAL-2} &
\textbf{SOV-7}\newline Security \& Compliance\newline \textit{10\% · min~SEAL-3} &
\textbf{SOV-8}\newline Environmental\newline \textit{5\% · min~SEAL-1} \\
\hline\hline
Q1 &
Hardware manufactured in EU or trusted jurisdictions? &
Non-proprietary APIs and open standards? &
Recognised EU certifications (ISO~27001, BSI~C5)? &
Energy-efficient infrastructure with low PUE? \\
\hline
Q2 &
Firmware of EU origin? &
Open-source licensed software stack? &
Compliant with GDPR, NIS2, and DORA? &
Circular economy practices for hardware? \\
\hline
Q3 &
Software stack architected and updated within EU? &
Full architectural transparency? &
All SOCs operating under EU jurisdiction? &
Transparent carbon and water disclosure? \\
\hline
Q4 &
No significant reliance on non-EU vendors? &
Free from non-EU HPC or processor dependency? &
Customer oversight of security monitoring? &
Renewable or low-carbon energy sourcing? \\
\hline
Q5 &
Full supply chain transparency and audit rights? &
--- &
EU-compliant breach reporting and patch management? &
--- \\
\hline
Q6 &
--- &
--- &
Independent EU security audit rights? &
--- \\
\hline
\end{tabularx}
\end{table}

The weights reflect the prioritisation established by the EC~DG~DIGIT framework. Supply Chain (SOV-5, 20\%) carries the highest weight, recognising that hardware and software dependencies represent the most structurally embedded sovereignty risk. Strategic Sovereignty (SOV-1) and Operational Sovereignty (SOV-4) jointly account for 30\%, reflecting the importance of EU-controlled governance and operational independence. Environmental performance (SOV-8, 5\%) carries the lowest weight, as it is not a sovereignty disqualifier under current EU regulation, though it is tracked for alignment with the EU Green Deal and the Corporate Sustainability Reporting Directive~(CSRD).

\subsection{Evidence Collection: The Y/P/N Verdict Protocol}
\label{subsec:ypn_verdicts}

For each of the 41 questions across the eight SOV objectives, a verdict is assigned per provider based on publicly available evidence. Three verdict values are used:

\begin{itemize}
  \item \textbf{Y (Yes)} — the provider fully and verifiably meets the criterion.
  \item \textbf{P (Partial)} — the provider partially meets the criterion, with documented limitations or caveats.
  \item \textbf{N (No)} — the provider does not meet the criterion, or no credible evidence of compliance was found.
\end{itemize}

Evidence sources include official provider documentation, annual reports, third-party certification records, press releases, and publicly disclosed regulatory filings. Each verdict is accompanied by a source URL and an explanatory note stored in a structured JSON survey file per provider. The evidence collection process was conducted systematically across all 18 providers using the same set of 41 questions, ensuring comparability.

Where a criterion cannot be assessed from public information alone, the conservative verdict~N is applied, consistent with the principle that sovereignty claims must be affirmatively demonstrable rather than assumed by default.%
\footnote{This is particularly relevant for SOV-5 (Supply Chain), where hardware firmware provenance is rarely disclosed publicly. Intel Management Engine and AMD PSP firmware, present in virtually all commercial x86 servers, are of non-EU origin, creating a structural ceiling that applies uniformly across all providers assessed.}

\subsection{SEAL Level Computation}
\label{subsec:seal_computation}

Each verdict is converted to a numeric score (Y~=~2, P~=~1, N~=~0). The raw score for a provider on a given SOV objective is the sum of its verdict scores across that objective's questions. This is mapped to a \textit{Sovereignty Assurance Level} (SEAL) on a five-point scale (0--4) using the thresholds defined in Table~\ref{tab:seal_levels}.

\begin{table}[htbp]
\centering
\caption{SEAL level thresholds and descriptions}
\label{tab:seal_levels}
\begin{tabular}{clp{9cm}}
\toprule
\textbf{SEAL} & \textbf{Name} & \textbf{Description} \\
\midrule
0 & No Sovereignty          & Service under exclusive control of non-EU third parties. \\
1 & Jurisdictional Sovereignty & EU law formally applies with limited practical enforceability. \\
2 & Data Sovereignty        & EU law applicable and enforceable; material non-EU dependencies remain. \\
3 & Digital Resilience      & EU law enforceable; EU actors exercise meaningful but not full influence. \\
4 & Full Digital Sovereignty & Technology and operations under complete EU control; no critical non-EU dependencies. \\
\bottomrule
\end{tabular}
\end{table}

The mapping from raw score to SEAL level uses fixed percentage thresholds applied to the maximum achievable score for each objective:

\begin{equation}
  \text{SEAL}_i =
  \begin{cases}
    0 & \text{if } p_i \leq 0.20 \\
    1 & \text{if } 0.20 < p_i \leq 0.40 \\
    2 & \text{if } 0.40 < p_i \leq 0.60 \\
    3 & \text{if } 0.60 < p_i \leq 0.80 \\
    4 & \text{if } p_i > 0.80
  \end{cases}
  \quad \text{where} \quad p_i = \frac{\text{score}_i}{\text{max\_score}_i}
  \label{eq:seal_mapping}
\end{equation}

The overall \textit{sovereignty score} for a provider is then computed as a weighted sum of normalised SEAL levels across all eight objectives:

\begin{equation}
  S_{\text{sov}} = \sum_{i=1}^{8} \frac{\text{SEAL}_i}{4} \times w_i \times 100 \quad (\%)
  \label{eq:sovscore}
\end{equation}

where $w_i$ is the weight of objective~$i$ (as listed in Table~\ref{tab:sov_objectives}) and dividing by~4 normalises each SEAL to the $[0,1]$ interval. The resulting score $S_{\text{sov}}$ lies in the range $[0, 100]$\%, where 100\% corresponds to a hypothetical provider achieving SEAL-4 on all eight objectives.

\subsection{Minimum SEAL Thresholds and Pass/Fail Determination}
\label{subsec:min_seal}

A sovereignty score alone does not determine whether a provider is fit for use in a regulated enterprise context. Each SOV objective carries a minimum acceptable SEAL level (see Table~\ref{tab:sov_objectives}), below which the provider is considered non-compliant on that objective regardless of its scores on other objectives. A provider \textit{passes} the assessment only if it meets or exceeds the minimum SEAL threshold on every objective simultaneously.

The minimum thresholds were set based on the regulatory environment applicable to enterprise cloud procurement in the EU. For instance, SOV-2 (Legal \& Jurisdictional) requires a minimum of SEAL-4, reflecting the binary nature of legal jurisdiction: a single residual channel through which a non-EU authority could compel data disclosure — such as the US CLOUD Act%
\footnote{The Clarifying Lawful Overseas Use of Data (CLOUD) Act (2018) allows US law enforcement to compel US-headquartered companies to produce data stored outside the United States, irrespective of the storage location. Its incompatibility with GDPR was confirmed by the CJEU in \textit{Schrems~II}~\cite{SchremsII_2020}.}
or the Chinese Cybersecurity Law — disqualifies a provider from processing personal data under GDPR~\cite{GDPR2016}. SOV-5 (Supply Chain) is set at SEAL-2, acknowledging the structural impossibility of full hardware sovereignty given the global x86 supply chain.

\subsection{Provider Coverage and Tool Implementation}
\label{subsec:sovscore_implementation}

The assessment was applied to 18 European cloud providers selected on the basis of EU headquarters, commercial availability, and relevance to enterprise workloads. For each provider, a structured JSON survey file was compiled containing the 41 verdicts with associated source URLs and explanatory notes.

The scoring framework was implemented as an interactive web application (SovScore) built with React~19 and FastAPI, allowing real-time exploration of results, adjustment of minimum SEAL thresholds, and comparison across providers. The application computes $S_{\text{sov}}$ from equation~(\ref{eq:sovscore}) in the browser on every state change, with all verdict data served statically from the JSON survey files.

\begin{figure}[htbp]
  \centering
  \fbox{\begin{minipage}{0.92\textwidth}
    \centering\vspace{2em}
    \textit{[SCREENSHOT REQUIRED: SovScore web interface showing the full provider $\times$ question grid,
    SEAL chips per SOV objective, and the Score column with sort order applied.
    Capture with all 18 providers visible and at least one provider highlighted in red (failing).]}
    \vspace{2em}
  \end{minipage}}
  \caption{SovScore web interface: provider assessment grid with SEAL levels and overall sovereignty scores.}
  \label{fig:sovscore_ui}
\end{figure}

\begin{figure}[htbp]
  \centering
  \fbox{\begin{minipage}{0.92\textwidth}
    \centering\vspace{2em}
    \textit{[SCREENSHOT REQUIRED: Hover tooltip on a SOV objective header (e.g.\ SOV-2 Legal \&
    Jurisdictional), showing the \textcircled{i} popover with the min-SEAL justification text.]}
    \vspace{2em}
  \end{minipage}}
  \caption{Min-SEAL justification tooltip on the SovScore grid header.}
  \label{fig:sovscore_tooltip}
\end{figure}


\subsection{Preliminary SovScore Findings}
\label{subsec:sovscore_preliminary}

The SovScore assessment surfaces a striking pattern across the 18 providers: the relationship between certification depth and sovereignty score is not linear. Providers holding the most extensive EU compliance portfolios — notably T-Systems (T~Cloud Public) with BSI~C5~Type~2, ISO~27001, SOC~1/2/3, TISAX, and GDPR Art.~28 DPAs — score near-perfectly on Security~\& Compliance (SOV-7) and Environmental (SOV-8) dimensions. Yet the same provider records the study's most complex sovereign weakness: a persistent partial score on Legal~\& Jurisdictional (SOV-2) and a failed verdict on Supply~Chain (SOV-5) due to the Huawei~FusionSphere dependency inherited from the original Open~Telekom~Cloud architecture. China's National Intelligence Law (2017) theoretically obligates Chinese companies to cooperate with state intelligence services, creating a residual extraterritorial risk that even the most rigorous German certification cannot fully eliminate.

A second observation is the universal failure of all EU providers on the Technology dimension (SOV-6) question regarding EU-origin high-performance computing hardware. No commercially available EU cloud provider deploys EU-manufactured server processors or GPU accelerators at scale; all depend on Intel or AMD x86 silicon (US origin) and NVIDIA accelerators (US origin). This is a structural constraint of the current semiconductor industry rather than a provider-specific deficiency, and it establishes a ceiling below which any EU cloud provider's SOV-6 score must fall regardless of governance choices.

A full cross-provider analysis of SovScore results, including ranked scores, SEAL-level breakdowns per objective, and pass/fail determinations against the minimum thresholds, is presented in Chapter~8.

% ─────────────────────────────────────────────────────────────────────────────
\section{Cloud Readiness Assessment Design}
\label{sec:readiness_methodology}
% ─────────────────────────────────────────────────────────────────────────────

\subsection{Rationale and Baseline Selection}
\label{subsec:readiness_rationale}

Sovereignty scoring measures whether a provider is legally and structurally safe to use under EU jurisdiction. It does not, however, measure whether a provider can actually replace the managed services that an enterprise already depends on. A cloud provider may achieve a high sovereignty score while lacking several services critical to a given workload, creating significant migration friction. The Cloud Readiness Assessment addresses this complementary question.

Amazon Web Services was selected as the 100\% capability baseline for three reasons. First, AWS holds the largest share of the global enterprise cloud market and is the de facto incumbent for the majority of European enterprises currently operating in a hyperscaler environment~\cite{armbrust2010view}. Second, AWS offers the most comprehensive publicly documented managed service portfolio, providing a stable reference taxonomy. Third, the use case driving this research — evaluating migration feasibility for Alps Alpine, operating on AWS — makes AWS the natural ``as-is'' baseline from which a transition to a European provider would depart.

The assessment is not a judgment on AWS's technical quality; rather, it uses AWS's service portfolio as a reference map of enterprise cloud capability requirements against which European alternatives are measured.

\subsection{Service Taxonomy: 11 Categories and 60 Services}
\label{subsec:service_taxonomy}

Sixty AWS managed services were selected as the evaluation rubric, organised into 11 functional categories. The selection was guided by two criteria: (1) services identified as required in the Alps Alpine IT infrastructure review, and (2) services appearing consistently in enterprise cloud reference architectures across the industry, as evidenced by the AWS Well-Architected Framework~\cite{AWS_WellArchitected_2023} and confirmed by enterprise adoption data in the Flexera State of the Cloud Report~\cite{Flexera_CloudReport_2024}. Services meeting both criteria are classified as \textit{required}; services meeting only the second are included as \textit{industry-standard additions} and are annotated accordingly in the tool.

Table~\ref{tab:readiness_services} lists all 60 services across the 13 assessment dimensions.
Services marked $^{\dagger}$ are industry-standard inclusions not specifically identified as
required by Alps Alpine but included to ensure full coverage of enterprise cloud reference
architectures~\cite{AWS_WellArchitected_2023, Flexera_CloudReport_2024}.

\begin{table}[htbp]
\centering
\caption{All 60 AWS services assessed across 13 readiness dimensions%
  \quad{\normalfont($^{\dagger}$ = industry-standard; not specifically required by Alps Alpine)}}
\label{tab:readiness_services}
\resizebox{\linewidth}{!}{%
\renewcommand{\arraystretch}{1.15}
\begin{tabularx}{1.6\linewidth}{XXXXXXX}
\toprule
\textbf{Scalability}        & \textbf{Performance}       & \textbf{Application (1)} &
\textbf{Compute (6)}        & \textbf{Storage (5)}       & \textbf{Database (7)}    &
\textbf{Networking (7)} \\
\midrule
\textit{Pre-assessed}       & \textit{Pre-assessed}      & Amplify                  &
Lambda                      & S3                         & RDS                      &
App.\ Load Balancer \\
                            &                            &                          &
ECS                         & S3 Glacier                 & DocumentDB               &
Net.\ Load Balancer \\
                            &                            &                          &
EKS                         & EBS                        & DynamoDB                 &
VPC \\
                            &                            &                          &
EC2                         & Backup                     & Timestream               &
Route~53 \\
                            &                            &                          &
Batch                       & EFS$^{\dagger}$            & Aurora$^{\dagger}$       &
API Gateway \\
                            &                            &                          &
Elastic Beanstalk$^{\dagger}$ &                          & ElastiCache$^{\dagger}$  &
CloudFront$^{\dagger}$ \\
                            &                            &                          &
                            &                            & Redshift$^{\dagger}$     &
Direct Connect$^{\dagger}$ \\
\midrule
\textbf{AI/ML (6)}          & \textbf{Security (12)}     & \textbf{Monitoring (4)}  &
\textbf{Dev/Deploy (5)}     & \textbf{IoT (2)}           & \textbf{Messaging (5)}   &
\\
\midrule
Rekognition                 & IAM                        & CloudWatch               &
CodeBuild                   & IoT Core                   & SQS                      & \\
SageMaker                   & Cognito                    & SNS                      &
CodeDeploy                  & IoT Device Mgmt.           & EventBridge              & \\
Bedrock                     & IAM Identity Center        & Systems Manager          &
CodePipeline                &                            & Step Functions           & \\
Glue                        & KMS                        & Cost Explorer            &
ECR                         &                            & SES                      & \\
EMR                         & Secrets Manager            &                          &
CloudFormation$^{\dagger}$  &                            & Kinesis$^{\dagger}$      & \\
Athena                      & Certificate Manager        &                          &
                            &                            &                          & \\
                            & GuardDuty                  &                          &
                            &                            &                          & \\
                            & Security Hub               &                          &
                            &                            &                          & \\
                            & Config                     &                          &
                            &                            &                          & \\
                            & Shield                     &                          &
                            &                            &                          & \\
                            & WAF                        &                          &
                            &                            &                          & \\
                            & CloudTrail$^{\dagger}$     &                          &
                            &                            &                          & \\
\bottomrule
\end{tabularx}
}% end resizebox
\end{table}

\subsection{Coverage Scoring Protocol}
\label{subsec:coverage_scoring}

For each of the 60 services, every EU provider under assessment receives a \textit{coverage verdict} indicating whether a functionally equivalent managed service exists in their portfolio:

\begin{itemize}
  \item \textbf{Y (Full)} — a managed equivalent service exists with comparable functionality and general availability status.
  \item \textbf{P (Partial)} — an equivalent exists but is in beta, feature-limited, or requires a workaround that increases operational overhead.
  \item \textbf{N (None)} — no managed equivalent is available; self-hosting on IaaS would be required, introducing additional operational burden.
\end{itemize}

This three-level coverage scheme was deliberately aligned with the Y/P/N verdict protocol used in the SovScore experiment, allowing a consistent scoring scale across both instruments. Where a partial equivalent exists, the EU service name or approach is recorded alongside the verdict to provide actionable migration guidance.

\subsection{Dimension Score Computation}
\label{subsec:dimension_score}

Coverage verdicts are converted to points (Y~=~2, P~=~1, N~=~0). The score for a provider on a given feature dimension $d$ is:

\begin{equation}
  R_d = \frac{\displaystyle\sum_{s \in d} \text{cov}(s)}{2 \times |d|} \times 100 \quad (\%)
  \label{eq:readiness_dim}
\end{equation}

where $|d|$ is the number of services in dimension~$d$ and $\text{cov}(s) \in \{0, 1, 2\}$ is the coverage point value for service~$s$. A score of 100\% indicates full managed-service parity with AWS on that dimension; 0\% indicates no managed equivalents exist and all services would require self-hosting.

Scalability and Performance are assessed separately as pre-scored dimensions (0--100) derived from published regional availability, auto-scaling capability documentation, and official SLA commitments, rather than from service-level verdicts.

The average readiness score across all 13 dimensions is used to rank providers and compute the \textit{average gap to AWS} as:

\begin{equation}
  \text{Gap} = \frac{1}{13} \sum_{d=1}^{13} \bigl(100 - R_d\bigr) \quad (\text{percentage points})
  \label{eq:readiness_gap}
\end{equation}

A smaller gap indicates higher overall feature parity with AWS and therefore lower expected migration friction.

\subsection{Tool Implementation and Visualisation}
\label{subsec:readiness_implementation}

The Cloud Readiness Assessment was implemented as the second module of the CloudSov web application, co-located with the SovScore tool. The frontend presents three complementary views of the assessment data:

\begin{enumerate}
  \item A \textbf{radar chart} showing each provider's coverage profile across all 13 dimensions simultaneously, with the AWS reference polygon as the outer boundary.
  \item An \textbf{area chart} showing dimension-by-dimension coverage with per-provider lines, enabling identification of specific capability gaps. Hovering over a data point reveals the coverage verdict for each service in that dimension.
  \item A \textbf{ranked gap table} listing providers by average gap to AWS, with expandable rows showing per-category service alternatives (EU equivalent name, or ``No alternative'' for N verdicts).
\end{enumerate}

\begin{figure}[htbp]
  \centering
  \fbox{\begin{minipage}{0.92\textwidth}
    \centering\vspace{2em}
    \textit{[SCREENSHOT REQUIRED: Cloud Readiness Assessment page showing the radar chart (left)
    and dimension breakdown area chart (right) with at least 4 providers visible, including AWS.
    Capture with the provider toggle buttons visible at the top.]}
    \vspace{2em}
  \end{minipage}}
  \caption{Cloud Readiness Assessment: coverage profile radar (left) and dimension breakdown area chart (right).}
  \label{fig:readiness_charts}
\end{figure}

\begin{figure}[htbp]
  \centering
  \fbox{\begin{minipage}{0.92\textwidth}
    \centering\vspace{2em}
    \textit{[SCREENSHOT REQUIRED: Gap summary table with one provider row expanded,
    showing the per-category service grid with EU alternative names and ``No alternative'' labels.
    Capture a provider that has a mix of Y, P, and N verdicts (e.g.\ Hetzner or Scaleway).]}
    \vspace{2em}
  \end{minipage}}
  \caption{Gap summary table with an expanded provider row showing per-service EU alternatives.}
  \label{fig:readiness_gap_table}
\end{figure}

The two visualisations are designed to serve complementary analytical purposes: the radar chart gives an immediate gestalt impression of a provider's overall readiness profile relative to AWS, while the area chart and gap table support drill-down analysis at the dimension and service level. Clicking any data point in the area chart pins a detailed service table below the chart, listing AWS service name, coverage verdict, EU equivalent, and source reference for every service in that dimension.


\subsection{Preliminary Readiness Findings}
\label{subsec:readiness_preliminary}

The readiness data reveals a pronounced stratification among EU providers that correlates closely with the depth and maturity of their managed-service portfolios. T-Systems OTC achieves the highest individual dimension scores among all EU providers, recording a scalability meta-score of~72 and a performance meta-score of~87 out of~100 — reflecting its wide managed-service catalogue including FunctionGraph (serverless), ModelArts (ML platform), IoTDA, DMS Kafka, and CodeArts. T Cloud Public, assessed separately as a distinct commercial product from the same parent, achieves a marginally higher scalability score of~74 and the same performance score of~87, reflecting its purpose-built sovereign infrastructure design and planned Netherlands capacity expansion. OVHcloud follows as the next-ranked EU provider (scalability~68, performance~83), driven by its breadth of managed databases and broad regional availability. Both providers score significantly higher than the infrastructure-centric providers at the lower end of the spectrum.

Hetzner, despite being one of Europe's most widely used cloud providers by customer count, records a scalability score of just~35 — reflecting its deliberate focus on virtual machines, block storage, and bare-metal servers without native managed services for serverless, workflow orchestration, or AI/ML. STACKIT, despite the substantial financial backing of the Schwarz~Group, scores similarly low on scalability~(32) owing to its 2021 launch date and the resulting managed-service portfolio still under active development. This pattern underlines the critical role of time-in-market: providers with deeper operational histories of cloud service delivery — particularly OVHcloud (cloud IaaS since~2012) and T-Systems~OTC (since~2016) — have broader catalogues than newer entrants.

Across all providers, the dimensions with the largest average gaps to the AWS baseline are AI/ML (Amazon~Bedrock has no equivalent among any EU provider), Database (no EU equivalent to DynamoDB, Amazon~Aurora, or Amazon~Timestream), and Dev/Deploy (Step~Functions has no EU managed equivalent). These structural gaps are consistent across providers of all sizes and ages, suggesting that they reflect the maturity of the underlying managed-service ecosystem rather than any individual provider's product strategy. Full findings with ranked provider scores and per-dimension gap analysis are discussed in Chapter~8.

% ─────────────────────────────────────────────────────────────────────────────
\section{Financial Analysis Methodology}
\label{sec:financial_methodology}
% ─────────────────────────────────────────────────────────────────────────────

\subsection{Rationale and Scope}
\label{subsec:financial_rationale}

The SovScore and Cloud Readiness Assessment instruments characterise \emph{what} EU providers offer and \emph{how} their jurisdictional posture compares to the AWS baseline. Both instruments, however, are silent on whether the capital underpinning EU providers is sufficient to close the observed readiness gap, and whether the \emph{source} of that capital has its own sovereignty implications. A third complementary analysis --- the Financial Consideration module --- addresses these questions directly. It asks: (1) at current investment velocity, when will EU providers reach AWS parity? (2) Is the capital already deployed proportionate to the readiness achieved? and (3) does the ownership and origin of that capital itself affect a provider's sovereignty score?

This section describes the Financial Consideration module of the CloudSov tool: how it was built, how its data were collected and classified, and how its five interactive visualisations should be read. The quantitative findings and their implications for the thesis research questions are discussed in Section~\ref{sec:financial_landscape}.

\subsection{Data Sources and Collection}
\label{subsec:financial_data_sources}

Financial data were assembled from four categories of publicly available sources:

\begin{enumerate}
  \item \textbf{Company filings and investor relations:} Annual reports, IPO prospectuses, and earnings releases from OVHcloud (Euronext Paris, since November~2021) and IONOS~SE (Frankfurt Stock Exchange, since February~2023). AWS revenue figures are drawn from Amazon's Form~10-K annual reports and converted to EUR at the 2024 annual average rate of \$1\,=\,€0.920.
  \item \textbf{Press releases and investor announcements:} Capital raises, PE/VC transactions, parent-company investment pledges, debt issuances, and government contract awards reported by each provider or their controlling group (Schwarz~Group, Iliad~Group, Deutsche~Telekom~AG).
  \item \textbf{IPCEI-CIS documentation:} The European Commission's December~2023 approval decision for the Important Project of Common European Interest on Next Generation Cloud Infrastructure and Services (IPCEI-CIS), including the official beneficiary list and published national funding allocations for Scaleway, OVHcloud, and Deutsche~Telekom.
  \item \textbf{EC sovereign cloud tender records:} The European Commission's April~2026 award notices for the €180M Sovereign Cloud Tender (Lots~1--3), which publicly disclosed winning consortia and per-lot contract values.
\end{enumerate}

Where a provider's financial data are not publicly disclosed --- most notably STACKIT (a wholly-owned Schwarz~Group internal unit with no separate reporting obligation), Scaleway (Iliad subsidiary), and T-Cloud~Public (Deutsche~Telekom subsidiary) --- revenue figures are derived from parent-group disclosures, analyst commentary, and cross-referencing with readiness capability indicators. These estimates are labelled as such in the tool interface. All monetary amounts are denominated in EUR; USD figures are converted using the rate above, applied uniformly across all years.

\subsection{Investment Milestone Classification}
\label{subsec:milestone_classification}

Each discrete capital event is recorded as an investment milestone with the following structured attributes: provider key, year, nominal EUR amount (where disclosed), event type, a short label, and a narrative description. Events with undisclosed amounts are retained in the milestone timeline for completeness but excluded from all quantitative aggregations. Event types are consolidated into six funding source categories for the capital structure analysis:

\begin{itemize}
  \item \textbf{Capex} --- direct infrastructure expenditure by the provider or its parent group (hardware procurement, data centre construction, regional expansion).
  \item \textbf{Debt} --- bank loans, revolving credit facilities, high-yield bonds, and EIB green loans.
  \item \textbf{Market Capital} --- IPO proceeds and private equity / venture capital injections.
  \item \textbf{Parent} --- discretionary capital commitments from the controlling parent entity (Schwarz~Group for STACKIT; Iliad for Scaleway) not separately categorised as Capex.
  \item \textbf{EU Programmes} --- IPCEI-CIS grant approvals and EC sovereign cloud tender contract awards.
  \item \textbf{Other} --- strategic acquisitions, infrastructure asset transactions, and government service contracts not classified under the above categories.
\end{itemize}

\subsection{Financial Consideration Module: Tool Implementation and Visualisations}
\label{subsec:financial_tool}

The Financial Consideration module is the third module of the CloudSov web application, accessible via the \textbf{Financial} tab in the top navigation bar. It presents five interactive charts, each addressing a distinct analytical question about the financial position and trajectory of EU sovereign cloud providers relative to AWS. Unlike the SovScore and Readiness modules, this module is read-only: there are no verdict inputs. All chart data are derived at backend startup from the structured milestone and revenue datasets described in Sections~\ref{subsec:financial_data_sources} and~\ref{subsec:milestone_classification}.

%─────────────────────────────────────────────────────────────────────────────
\subsubsection{Readiness Gap Closure — EU Providers vs AWS}
\label{subsubsec:chart_gap_closure}
%─────────────────────────────────────────────────────────────────────────────

Figure~\ref{fig:financial_row1_left} shows each EU provider's historical readiness trajectory as a solid coloured line from its cloud launch year to the current score, followed by a dashed projection at the same velocity until the score reaches 100 (AWS parity). The underlying model is linear: for a provider with cloud launch year $y_0$ and average readiness score $s$ at Q2~2026, the annualised velocity and projected parity year are defined as:

\begin{align}
  v &= \frac{s}{2026 - y_0} \quad [\text{pts/yr}] \label{eq:velocity} \\[4pt]
  t_{\text{parity}} &\approx 2026 + \frac{100 - s}{v} \label{eq:parity}
\end{align}

\noindent where $v$ is the annualised readiness velocity and $t_{\text{parity}}$ is the projected year at which the provider reaches score\,=\,100. Estimated parity years derived from Equations~\ref{eq:velocity}--\ref{eq:parity} are annotated on each trajectory. Two vertical reference lines are overlaid: a grey dashed marker at 2026 labelled \emph{Now}, indicating the analysis cut-off, and an amber dashed marker at 2030 labelled \emph{EC Digital Decade}, corresponding to the headline cloud adoption target of the EU Digital Decade Policy Programme~\cite{EC_DigitalDecade_2022} (see Section~\ref{sec:digital_decade}). The chart addresses the question: \emph{at current investment velocity, when will each EU provider reach AWS readiness parity, and which providers are on track to do so within the Digital Decade horizon?}

\begin{figure}[htbp]
  \centering
  \fbox{\begin{minipage}{0.92\textwidth}
    \centering\vspace{2em}
    \textit{[SCREENSHOT REQUIRED: Financial Consideration page --- Readiness Gap Closure projection chart showing solid historical lines and dashed projected trajectories for all five EU providers, with AWS parity line at score\,=\,100, annotated parity year labels, the grey ``Now'' marker at 2026, and the amber ``EC Digital Decade'' marker at 2030.]}
    \vspace{2em}
  \end{minipage}}
  \caption{Readiness gap closure: historical readiness velocities (solid) projected forward (dashed) to AWS parity (score\,=\,100). The amber vertical line marks the EC Digital Decade 2030 target~\cite{EC_DigitalDecade_2022}.}
  \label{fig:financial_row1_left}
\end{figure}

%─────────────────────────────────────────────────────────────────────────────
\subsubsection{Annual Revenue Scale — EU Providers vs AWS}
\label{subsubsec:chart_revenue}
%─────────────────────────────────────────────────────────────────────────────

Figure~\ref{fig:financial_row1_right} plots annual cloud-segment revenue on a log-scale Y-axis for AWS and the five EU providers across 2019--2024. A logarithmic scale is necessary because plotting AWS and the EU cohort on a linear axis would render the EU provider lines indistinguishable from zero: the multi-order-of-magnitude difference between AWS (€92B in 2024) and the highest-revenue EU provider (IONOS, €1.3B) is only legible on a log scale. The chart addresses the question: \emph{how large is the revenue gap between AWS and EU providers, and is it narrowing or widening?}

\begin{figure}[htbp]
  \centering
  \fbox{\begin{minipage}{0.92\textwidth}
    \centering\vspace{2em}
    \textit{[SCREENSHOT REQUIRED: Financial Consideration page --- Revenue Scale Comparison (log Y-axis) showing AWS and the five EU providers as coloured lines, 2019--2024. The AWS line should be clearly visible at the top of the log scale; EU provider lines clustered at the lower end.]}
    \vspace{2em}
  \end{minipage}}
  \caption{Annual cloud-segment revenue (EUR~M, log scale): AWS vs the five EU providers, 2019--2024. A log scale is required to render both AWS and EU providers on a common axis.}
  \label{fig:financial_row1_right}
\end{figure}

%─────────────────────────────────────────────────────────────────────────────
\subsubsection{Investment and Funding Milestone Timeline}
\label{subsubsec:chart_milestones}
%─────────────────────────────────────────────────────────────────────────────

Figure~\ref{fig:financial_row2} presents a full-width bubble chart that places each recorded investment milestone against calendar year on the X-axis. Bubble area is proportional to the committed amount (EUR); colour encodes the provider. Hovering over a bubble reveals the event label, year, amount, and event type. Milestones with undisclosed amounts are displayed at a fixed minimum size with a dashed outline, distinguishing them from confirmed capital events. AWS milestones are included for scale reference. The chart addresses the question: \emph{what discrete capital events have shaped each provider's trajectory, and how do they compare in scale to AWS's infrastructure commitments?}

\begin{figure}[htbp]
  \centering
  \fbox{\begin{minipage}{0.92\textwidth}
    \centering\vspace{2em}
    \textit{[SCREENSHOT REQUIRED: Financial Consideration page --- Investment and Funding Milestone Timeline bubble chart. Capture with STACKIT's €11B Bitterfeld-Wolfen bubble clearly visible (the largest bubble in the chart) alongside the AWS reference bubbles. All provider colours should be visible in the legend. Dashed-outline bubbles indicate undisclosed amounts.]}
    \vspace{2em}
  \end{minipage}}
  \caption{Investment and funding milestones for AWS and the five EU providers, 2006--2026. Bubble area is proportional to the committed capital (EUR); dashed outlines indicate undisclosed amounts.}
  \label{fig:financial_row2}
\end{figure}

%─────────────────────────────────────────────────────────────────────────────
\subsubsection{Investment Efficiency — Readiness Points per Euro Invested}
\label{subsubsec:chart_efficiency}
%─────────────────────────────────────────────────────────────────────────────

Figure~\ref{fig:financial_row3_left} is an investment efficiency scatter in which each EU provider is plotted with cumulative disclosed capital on the log-scale X-axis (€M) and current average readiness score on the Y-axis. The efficiency ratio $\eta$ plotted beneath each provider label is defined as:

\begin{equation}
  \eta = \frac{s}{C_{\text{disclosed}}} \quad [\text{pts/\texteuro{}B}] \label{eq:efficiency}
\end{equation}

\noindent where $s$ is the provider's average readiness score and $C_{\text{disclosed}}$ is the cumulative disclosed capital in billions of euros. Only capital events with a publicly confirmed nominal amount contribute to $C_{\text{disclosed}}$; milestones with undisclosed amounts are excluded from the denominator. Each data point displays the provider name above and $\eta$ below. Providers positioned towards the upper-left --- high readiness relative to low disclosed capital --- represent the most capital-efficient paths to sovereignty. The chart addresses the question: \emph{for every euro disclosed as invested, how much measurable cloud readiness has each provider achieved?}

\begin{figure}[htbp]
  \centering
  \fbox{\begin{minipage}{0.92\textwidth}
    \centering\vspace{2em}
    \textit{[SCREENSHOT REQUIRED: Financial Consideration page --- Investment Efficiency scatter chart. Each EU provider plotted as a labelled dot with provider name above and pts/€B ratio below. Log X-axis (cumulative disclosed capital, €M), linear Y-axis (readiness score 0--100).]}
    \vspace{2em}
  \end{minipage}}
  \caption{Investment efficiency scatter: current average readiness score vs cumulative disclosed capital (log scale, €M). The efficiency ratio $\eta$ (readiness points per €B) is shown beneath each provider label.}
  \label{fig:financial_row3_left}
\end{figure}

%─────────────────────────────────────────────────────────────────────────────
\subsubsection{Capital Structure by Funding Source}
\label{subsubsec:chart_capital_structure}
%─────────────────────────────────────────────────────────────────────────────

Figure~\ref{fig:financial_row3_right} is a stacked bar chart decomposing each provider's total disclosed capital into the six funding source categories defined in Section~\ref{subsec:milestone_classification}. Each bar segment is colour-coded by category (Capex, Debt, Market Capital, Parent, EU Programmes, Other), and the chart is sorted to make the dominant funding model of each provider immediately legible. The chart addresses the question: \emph{where does each provider's capital come from, and does that origin carry its own sovereignty implications under the SovScore framework?}

\begin{figure}[htbp]
  \centering
  \fbox{\begin{minipage}{0.92\textwidth}
    \centering\vspace{2em}
    \textit{[SCREENSHOT REQUIRED: Financial Consideration page --- Capital Structure stacked bar chart. Each provider shown as a stacked bar with segments colour-coded by funding category (Capex, Debt, Market Capital, Parent, EU Programmes, Other). STACKIT's bar should dominate the chart due to the €11B parent commitment; OVHcloud and IONOS should show Market Capital segments.]}
    \vspace{2em}
  \end{minipage}}
  \caption{Capital structure by funding source: total disclosed capital decomposed by category for each EU provider (€M). Undisclosed events are excluded.}
  \label{fig:financial_row3_right}
\end{figure}

\subsection{Preliminary Financial Findings}
\label{subsec:financial_preliminary}

The financial analysis surfaces three immediate observations that contextualise the SovScore and Readiness results established in preceding sections.

First, the revenue gap between AWS and the EU cohort is not merely large --- it is structurally widening. AWS's sustained ~17\% annual CAGR means the absolute revenue gap grows by €13--16B per year even as EU providers achieve their own growth. No EU provider in the study is on a trajectory to close this gap through organic revenue growth alone in a policy-relevant timeframe.

Second, the gap-closure projection reveals a two-tier EU cohort. T-Systems (OTC / T-Cloud Public) and STACKIT are on trajectories that place AWS readiness parity within reach by approximately 2029--2031, driven by a combination of operational maturity (T-Systems) and accelerating capital deployment (STACKIT). OVHcloud, Scaleway, and IONOS, by contrast, face parity timelines extending into the late 2030s and 2040s under the current investment velocity.

Third, and most directly relevant to the thesis research question on sovereignty, the capital structure decomposition reveals a clean divide: providers financed through EU parent groups or EU institutional channels (STACKIT, Scaleway, T-Systems) achieve full \texttt{Y} verdicts on the SovScore questions that directly score capital origin (SOV-1~Q3, SOV-2~Q2). Providers with US private equity or public market capital exposure (OVHcloud via KKR and General Atlantic; IONOS via Warburg~Pincus) receive \texttt{P} verdicts on these same questions. Capital structure is not merely a financial characteristic --- it is a direct sovereignty determinant within the scoring framework used in this thesis. A full discussion of these findings, including the causal mechanism through which capital origin affects sovereignty scoring, is presented in Section~\ref{sec:financial_landscape}.

\subsection{Limitations}
\label{subsec:financial_limitations}

The financial analysis is subject to several important limitations. First, only \emph{disclosed} capital events contribute to quantitative aggregations; providers with limited public reporting obligations (STACKIT, Scaleway, T-Cloud~Public) are understated in the capital efficiency and capital structure analyses. Second, the revenue series for non-listed providers are management estimates subject to revision. Third, investment milestones represent \emph{announced} or \emph{committed} amounts, which may differ from capital actually deployed --- most notably STACKIT's €11B Bitterfeld-Wolfen campus pledge, where Phase~1 construction is targeted for completion in~2027 and the associated managed services have not yet been delivered. Fourth, the uniform EUR/USD conversion rate does not account for multi-year exchange-rate fluctuation. These limitations are acknowledged in the discussion in Section~\ref{sec:financial_landscape}.

% ─────────────────────────────────────────────────────────────────────────────
\section{Benchmark Pipeline Design}
\label{sec:benchmark_methodology}
% ─────────────────────────────────────────────────────────────────────────────

...

\subsection{Compute (sysbench)}
\subsection{Storage (fio)}
\subsection{Network (iperf3)}
\subsection{AI/ML Inference (Odin)}


% ─────────────────────────────────────────────────────────────────────────────
\section{Cost Evaluation Methodology (TCO)}
\label{sec:tco_methodology}
% ─────────────────────────────────────────────────────────────────────────────

...


% ─────────────────────────────────────────────────────────────────────────────
\section{Deployment Experiment Design (Phase 3)}
\label{sec:deployment_methodology}
% ─────────────────────────────────────────────────────────────────────────────

...


% ─────────────────────────────────────────────────────────────────────────────
\section{Evaluation Criteria and Metrics Summary}
\label{sec:evaluation_criteria}
% ─────────────────────────────────────────────────────────────────────────────

...
